%
% tikztemplate.tex -- Template für standalone TIKZ Bilder
%
% (c) 2019 Prof Dr Andreas Müller, Hochschule Rapperswil
%
\documentclass[tikz,12pt]{standalone}
\usepackage{amsmath}
\usepackage{times}
\usepackage{txfonts}
\usepackage{pgfplots}
\usepackage{csvsimple}
\usetikzlibrary{arrows,intersections,math,calc}
\definecolor{darkred}{rgb}{0.8,0,0}
\definecolor{darkgreen}{rgb}{0,0.6,0}
\begin{document}
\def\skala{1}
\begin{tikzpicture}[>=latex,thick,scale=\skala]

\def\s{2}
\pgfmathparse{1+ln(2)}
\xdef\C{\pgfmathresult}
\xdef\D{-1}
\begin{scope}
	\clip (2,-4.1) rectangle (10.2,4.1);
	\fill[color=blue!10] (2,-4.1) rectangle (10.2,4.1);
	\fill[color=orange!20]
		plot[domain=1:2,samples=10]
			({\s*\x},{\s*(0.5*\x*\x+\D)})
		--
		plot[domain=2:1,samples=10]
			({2*\x},{\s*(-ln(\x)+\C)})
		-- cycle;
	\foreach \x in {2,4,...,11}{
		\draw[color=white] (\x,-4.1) -- (\x,4.1);
	}
	\foreach \y in {-4,-2,...,4}{
		\draw[color=white] (1,\y) -- (10.2,\y);
	}
\end{scope}
\draw[color=orange,line width=2pt] (2,{\s*(0.5+\D)}) -- (2,{\s*\C});
\node[color=orange!80!black] at (2,{\s*0.5*(0.5+\D+\C)}) [left] {$G$};
\node[color=blue] at (5,0.5) {$\Omega$};
\begin{scope}
	\clip (2,-4.1) rectangle (10.2,4.1);
	\foreach \D in {-15,-14.5,...,2}{
		\draw[color=darkred,line width=1.4pt] 
			plot[domain=1:5.1,samples=100]
				({2*\x},{\s*(0.5*\x*\x+\D)});
	}
	\foreach \C in {-2.5,-2,...,2.5}{
		\draw[color=darkgreen,line width=1.4pt] 
			plot[domain=1:5.1,samples=100]
				({2*\x},{\s*(-ln(\x)+\C)});
	}
	\end{scope}
\fill[color=blue] (4,2) circle[radius=0.08];
\node[color=blue] at (4,2) [above right] {$P$};
\draw[->] (-0.1,0) -- (10.7,0) coordinate[label={$x$}];
\draw[->] (0,-4.1) -- (0,4.4) coordinate[label={right:$y$}];
\foreach \y in {-4,-2,2,4}{
	\draw (-0.05,\y) -- (0.05,\y);
}
\node at (-0.05,-4) [left] {$-2\mathstrut$};
\node at (-0.05,-2) [left] {$-1\mathstrut$};
\node at (-0.05,0) [left] {$0\mathstrut$};
\node at (-0.05,2) [left] {$1\mathstrut$};
\node at (-0.05,4) [left] {$2\mathstrut$};
\foreach \x in {1,...,5}{
	\draw ({\s*\x},-0.05) -- ({\s*\x},0.05);
	\node at ({\s*\x},0) [below left] {$\x\mathstrut$};
}

\end{tikzpicture}
\end{document}

